Pragmatic sentiment in Vietnamese text is not exhausted by polarity. Sarcasm, irony, idiomaticity, implicitness, code-switching, and mocking can change how a reader should interpret an utterance even when its surface sentiment is positive or negative. We introduce ViPragSent, an XLM-R-large encoder with six distinct binary pragmatic heads, separate three-way polarity and seven-way emotion auxiliaries, and a training-only teacher-forced rationale objective. The model uses a shared first-nonpadding-token representation and homoscedastic uncertainty weighting to couple these tasks while keeping their semantic roles separate. On a matched 2,000-example ID/gold test cohort, ViPragSent reaches 93.67 $\pm$ 0.19 macro-pragmatic F1 versus 92.93 $\pm$ 0.14 for the standard XLM-R baseline, an observed +0.73 percentage-point difference. Its mean is also highest on all six pragmatic heads and on the macro average among five complete three-seed baseline families. Follow-up analyses show that the full objective is an operating point rather than a uniform winner: Q2 reports accuracy, calibration, and cost trade-offs; Q3 reports a non-monotonic positive-label budget curve with contextual baselines; and Q4 reports target ECE 0.0386 $\pm$ 0.0069 versus Vistral 0.0345 $\pm$ 0.0041. These results define what the implemented model contributes while keeping rationale supervision separate from inference.
